\section{Data, stalls and the estimator}
\label{sec:data}

\subsection{The CS1 data and the holdout}

The confirmatory data are the 2019 edition of a public keystroke dataset from an introductory programming course at a
US university, in which students wrote Python in a browser-based editor and every edit was logged in the ProgSnap2
format \cite{edwards2023cs1}. A writer is a student and an assignment is one unit; within an assignment, events more than
30 minutes apart start a new work block. Before any analysis, 146 of the 487 writers were drawn at random as the holdout
(Appendix~\ref{app:prereg}). Of these, 145 had at least one assignment (one writer had a single checkpoint row and no key
press), giving 438 assignments; 140 had at least 10 minutes of observation, and 139 had at least one stall paired with
a control. Of the 341 seen writers, 338 had an assignment, 319 had at least 10 minutes of observation and 312 had a stall
paired with a control; these are the seen writers of the descriptions [D] and of the calibration.

The data record edit events rather than key presses, so the input to the estimator was built by a mapping fixed before
the study (Appendix~\ref{app:definitions}). An insertion of one character with a recorded key became a key press of that
character's kind, a deletion with a recorded key became a backspace, and an undo or redo with a key became a non-text
key press. Edits without a recorded key, pastes, cuts and drags did not become key presses but were kept as events, and so were
runs of the program and switches between tasks; only runs, pastes and switches define the event stratum. Because the data have no key release, the release was placed 100~ms
after each press but no later than 1~ms before the next press, so the estimator's flight time is the inter-keystroke
interval minus 100~ms when the interval exceeds 100~ms, 1~ms for shorter intervals and 0 for simultaneous presses, a
definition that differs from that of recorded releases in writing data. Table~\ref{tab:mapping} (Appendix~\ref{app:definitions}) shows how the edits of the holdout
writers were distributed. Every insertion with a recorded key was a single character. Insertions of several characters were rare
(1.4\% of edit rows) but carried 53.7\% of the inserted characters, mostly through pastes; none of them reached the
estimator as key presses.


\subsection{Pauses and stalls}

A \emph{pause} is the interval from one character input (a letter, digit, symbol, space, newline or tab) to the next
character input in the same assignment, so it can contain key presses that are not characters, such as deletions. A
stall is defined by a rule carried over from an earlier preregistered definition (Appendix~\ref{app:definitions}). Let
condition A hold at time $t$ when, in the window $(t - 30\,\mathrm{s}, t]$, the number of insertion presses minus the
number of deletion presses is less than one. When condition A rises from false to true at time $s$ and stays true until
$t_0 = s + \theta$, an onset is raised at $t_0$. Here $\theta$ is the 0.99 quantile of the writer's inter-keystroke
intervals so far in the assignment, defined once at least 100 intervals exist; it personalises the moment at which the
stall starts. We write $u$ for the last character input at or before $t_0$. An onset becomes a \emph{stall} when the
pause that starts at $u$ lasts at least 30~s, $u$ is earlier than $s$, and condition A is false at $u$ (the writer had
net production in the preceding 30~s). A stall is therefore not simply 30~s without character input: condition A must
first become true and then stay true for $\theta$, so the onset lies well after $u$. In the holdout's main stratum
[E-H], the onset came a median of 48.5~s after $u$ (10th and 90th percentiles 35.2 and 77.6~s), and the next character
input a median of 97.2~s after $u$ (78.6\% of stalls lasted 60~s or more; a few very long stops raise the mean to 247~s and
428~s; Appendix~\ref{app:definitions} gives the edges of observation). The 0.99 quantile was inherited from the earlier definition and was not chosen
in this study; the 30~s length of a stall and the 2~s lower limit of the controls below are fixed.

What happens inside a pause is not used to match stalls and controls. In the holdout's main stratum
(Section~\ref{sec:design}), 47\% of stalls contained a run of the program, a paste or a switch inside the pause,
against 6\% of controls; 52\% of stalls and 69\% of controls contained no key press of any kind, and the 90th percentile
of key presses inside the pause was 13 for stalls and 5 for controls.

\subsection{The estimator and its two signals}

The estimator reads the key presses and releases and, every second, computes the five quantities named in
Section~\ref{sec:intro} over the preceding 30~s. Each quantity is mapped to $[0, 1]$ against a reference value and combined into two latent axes, a friction
axis ($0.30$ correction ratio $+ 0.25$ pauses after deletion $+ 0.25$ flight time $+ 0.20$ number of pauses) and an
engagement axis, whose exponentially smoothed values are quantised into the observations of a three-state hidden Markov
model updated by the forward algorithm. The designer named the states Flow, Incubation and Stuck; we use only the
probability of the third and do not interpret the names. All coefficients were set by hand (Section~\ref{sec:intro}) and are listed in
Appendix~\ref{app:coefficients}; the version used is fixed (Appendix~\ref{app:definitions}). A replay tool re-reads a recorded stream and writes one
estimate row per second together with the stall onsets.

We use two signals, each taken as the change over 30~s: the \emph{change in the Stuck-state probability}, the
probability at time $t$ minus the probability 30~s earlier, and the \emph{change in the smoothed friction axis}, defined
in the same way. We write $\Delta p_{\mathrm{stuck}}$ and $\Delta x$ for them in tables. The levels at $t$ are reported
as descriptions. The Stuck-state probability has a floor: in the holdout as a whole, 40.5\% of the 2,602,660 estimate
rows were below 0.02, and the share per writer ranged from 26\% to 62\% between the 10th and 90th percentiles. Both
signals nevertheless moved on this input. The standard deviation of the 30-second change over all estimate rows was
0.270 for the Stuck-state probability and 0.187 for the friction axis, and within stalls both rose after $u$: the mean
Stuck-state probability went from 0.135 at $u$ to 0.193 at $u + 20$~s, and the mean friction axis from 0.291 to 0.412
(exploration on the holdout; Appendix~\ref{app:sensitivity}).

\subsection{The essay-writing and co-writing data}

KLiCKe is a corpus of argumentative essays written by 4,992 students in one session each, with recorded press and
release times \cite{tian2025klicke}. Sessions with fewer than 200 key presses are excluded, the boundary type is whether
the key press at $u$ is a space, newline or sentence punctuation, and no event stratum is used. We had analysed it
before this study, so its results are descriptions on already seen data. CoAuthor, a dataset of writing with suggestions
from a language model \cite{lee2022coauthor}, appears only in the exploration (Appendix~\ref{app:explore}).
